%! TeX program = lualatex
%tags! energia fisica cinetica 
\documentclass[letterpaper]{article}
\usepackage[margin=1in]{geometry}
\usepackage{amsmath}
\usepackage{amssymb}
\usepackage{circuitikz}
\usepackage[no-math]{fontspec}
\usepackage[]{gruvboxpalette}
\usepackage{hyperref}
\usepackage{newtxsf}
\usepackage[explicit]{titlesec}
\usepackage{tikz}
\usepackage{varwidth}
\usetikzlibrary{calc}
\usetikzlibrary{positioning}
\usetikzlibrary{arrows.meta}
\usepackage[most]{tcolorbox}
\usepackage{tabularray}
\DefTblrTemplate{firsthead, middlehead,lasthead}{default}{}
\DefTblrTemplate{capcont}{default}{}
\DefTblrTemplate{contfoot-text}{normal}{} \SetTblrTemplate{contfoot-text}{normal} \DefTblrTemplate{conthead-text}{normal}{} \SetTblrTemplate{conthead-text}{normal}
\UseTblrLibrary{counter}
\hypersetup{
  colorlinks  = true,
  urlcolor    = Blue,
  linkcolor   = Blue,
  citecolor   = Blue
}
\usepackage[most]{tcolorbox}

\setmainfont{JetBrainsMono-Regular}[
Path           = /home/snouflake/.fonts/ ,
Extension      = .ttf ,
BoldFont       = JetBrainsMono-Bold ,
ItalicFont     = JetBrainsMono-Italic ,
BoldItalicFont = JetBrainsMono-BoldItalic,
]  

\newtcolorbox{defbox}[3][]{%
  colback=blue!30!background,
  coltitle=blue!15!black,
  coltext=font,
  title filled=false,
	enhanced,
  detach title,
  tile,
  before upper={\tcbtitle\medskip\\},
  borderline west={2mm}{0pt}{blue},
  % attach boxed title to top center={yshift=-2mm},
  leftrule=2mm,
  toprule=0mm,
  bottomrule=0mm,
  rightrule=0mm,
  arc=0mm,
	title={Definición:~#2},
	#1
}

\setlength\parindent{0pt}

\usepackage[italic]{mathastext}

\def \T{Física 1}
\def \S{Energía}

\begin{document}
\begin{tikzpicture}[inner sep=0pt,color=font]
  \node[anchor=west,align=left,line width=0pt] 
    (title) at (0,0) {\huge\bfseries\noindent\T}
    ;
  \node[anchor=north west,align=left,line width=0pt] 
    (subtitle) at (title.south west) {\Large\bfseries\S\strut}
    ;
\end{tikzpicture}


\begin{longtblr}{
    colspec={@{}Q[4cm,cmd=\textbf,h] >{\begin{minipage}{\linewidth}}X<{\end{minipage}}@{}},
    rowsep={7pt},
  }
  Energía cinética
  &
  \[
    E_c = K = \frac{1}{2}\, m \left|v\right|^2 \hspace{10pt} \text{J\,[oules]}
  \]
  * Normalmente $\left|v\right|$, la rapidez, se escribe como $v$, pero tanto $E_c$, y $v$ en este caso, son escalares.

  \\
  Energía potencial gravitacional
  &
  \begin{equation}
    E_p = U = m g h
    \label{eq:ep}
  \end{equation}
  *donde h es distancia hacia abajo ($-y$), no arriba

  \begin{equation}
    \Delta U = -W \Rightarrow F(x) = -\nabla U(r)
  \end{equation}

  \\
  Energía mecánica total
  &
  \begin{equation}
    E_{mt} = E_c + E_c
    \label{eq:emt}
  \end{equation}
  \\
  Principio de la conservación
  &
  \begin{equation}
    E_{mt\,i} = E_{mt\,f}; \Delta K = - \Delta U
    \label{eq:emtemtf}
  \end{equation}
  \medskip

  La fuerza, por lo tanto, es conservativa ($\nabla \times \bar{F} = 0$), por lo cual también tiene una función de potencial.
  \\
  Trabajo
  &
  \begin{equation}
    W = \int \bar{F} \cdot d\bar{r}
    \label{eq:wdef}
  \end{equation}
  Si la energía está asociada a un potencial entonces $W$ no depende de la trayectoria.

  La fricción y la resistencia por gravedad no son conservativas
   \\
   Teorema del trabajo energía
  &
  \begin{equation}
    W = \Delta K = K_f - K_i
    \label{eq:teoremata}
  \end{equation}
  \\
  Plano inclinado
  &
  Sí
  \\
  Ley de Hooke
  &
  \begin{equation}
    F = -kx
    \label{eq:hooke}
  \end{equation}
  \\
  Trabajo en Hooke
  &
  \begin{equation}
    W = \int_{x_i}^{x_f} -kx dx
  \end{equation}
  \\
  Potencia???
  &~
  \\
  Momento
  &
  \begin{equation}
    \vec{p} = m\vec{v}
  \end{equation}

  Según Newton:
  \begin{equation}
    \vec{F}_{\text{neta}} = \frac{d\vec{p}}{dt} = \frac{d}{dt} (m\vec{v})
  \end{equation}

  \bigskip
  \textbf{[Inter]cambio de momento entre dos partículas}
  \begin{align*}
    \Delta p_1 &= -\Delta p-2 \\
    m_1 \Delta v_1 &= - m_2 \Delta v_2 \\
  \end{align*}

  \textbf{Cambio de momento}
  \begin{equation}
    \Delta p = \vec{p_f} - \vec{p_i}
  \end{equation}

  El momento total en un sistema aislado es constante.
  \\
  Centro de masa
  &
  \begin{equation}
    cdm_{x} = \frac{\sum^{n}_{i}m_i x_i }{\sum^{n}_{i}m_i}x
    \label{eq:cdmx}
  \end{equation}

  \begin{align}
    M \vec{r}_{cdm} = m_1 \vec{r}_1 + m_2 \vec{r}_2 + \dots \\
    M \vec{v}_{cdm} = m_1 \vec{v}_1 + m_2 \vec{v}_2 + \dots \\
    M \vec{a}_{cdm} = m_1 \vec{a}_1 + m_2 \vec{a}_2 + \dots
    \label{eq:mrcdm}
  \end{align}
  \\
  Colisiones
  &
  \begin{align}
    \vec{F} = \frac{d\vec{p}}{dt}
  \end{align}
  
  Impulso
  \begin{equation}
    \vec{J} = \int^{t_f}_{t_i} \vec{F}(t)\,dt = p_f - p_t \frac{kg\,m}{s}
  \end{equation}
\end{longtblr}


\end{document}




